\section{Independence Testing}

In this lecture, we will consider the following questions. How can we test whether\dots
\begin{itemize}
  \item \dots two random variables are independent?
  \item \dots two random variables are conditionally independent given a third random variable?
  \item \dots a random variable is independent of a non-random variable?
  \item \dots a conditional random variable is conditionally independent of a conditional random variable, given another conditional random variable?
\end{itemize}
We will consider these questions only for the special case of finite categorical variables, i.e., 
variables that take values in finite spaces. In particular, we will discuss a test known as the $G$ test.
This has been defined in the literature for random variables, but we will extend it here to a general case
involving conditional random variables (with ``purely'' random and ``purely'' non-random variables as special cases).
We will state conditions under which the tests are asymptotically valid and consistent.
%We will discuss frequentist tests, and if time allows, a Bayesian test as well.
%We will briefly mention one other popular test, the partial correlation test.

\subsection{Marginal Independence for Categorical Random Variables}

Consider two categorical random variables $X,Y$ taking values in finite spaces $\C{X}$ and $\C{Y}$,
respectively, with $2 \le |\C{X}| < \infty$ and $2 \le |\C{Y}| < \infty$, and joint distribution $P(X,Y)$.
We can represent the density in a table (assuming $\C{X} = \{1,\dots,k\}$ and $\C{Y} = \{1,\dots,l\}$):
\begin{center}\begin{tabular}{c|cccc|c}
  & $Y=1$         & $Y=2$ & $\dots$ & $Y=l$ \\
  \hline
  $X=1$ & $\theta_{11}$      & $\theta_{12}$      & $\dots$ & $\theta_{1l}$ & $\theta_{1+}$\\
  $X=2$ & $\theta_{21}$      & $\theta_{22}$      & $\dots$ & $\theta_{2l}$ & $\theta_{2+}$\\
  $\vdots$ & $\vdots$     & $\vdots$       & $\ddots$ & $\vdots$  & $\vdots$ \\
  $X=k$ & $\theta_{k1}$      & $\theta_{k2}$      & $\dots$ & $\theta_{kl}$ & $\theta_{k+}$ \\
  \hline
  & $\theta_{+ 1}$ & $\theta_{+ 2}$ & $\dots$ & $\theta_{+ l}$ & $\theta_{++}=1$\\
\end{tabular}\end{center}
where we introduced the parameter $\theta \in \CX_\Theta$ by setting $\theta_{xy} = P(X=x,Y=y)$ for
$x \in \C{X}, y \in \C{Y}$.
We introduce here the convention that a ``$+$'' index denotes summation over that index, i.e.,
$$\theta_{+ y} := \sum_{x\in\C{X}} \theta_{xy}, \qquad \theta_{x+} := \sum_{y\in\C{Y}} \theta_{xy}, \qquad \theta_{++} := \sum_{x\in\C{X}}\sum_{y\in\C{Y}} \theta_{xy}.$$
For the parameter space we take the $(|\C{X}||\C{Y}|-1)$-dimensional simplex:
$$\CX_\Theta := \{ \theta \in \prod_{(x,y)\in\C{X}\times\C{Y}} [0,1] : \theta_{+ +} = 1 \}.$$
With Remark~\ref{ci-random-variables}, we get:
$$X \Indep_{P(X,Y)} Y \iff P(X,Y) = P(X) \otimes P(Y),$$
where $P(X)$ and $P(Y)$ are the marginal distributions of $P(X,Y)$.
In the discrete case we consider here, this holds if and only if
$$\forall x \in \C{X}, y \in \C{Y}: \theta_{xy} = \theta_{x+} \theta_{+ y}.$$
The parameters satisfying this constraint form the allowed parameters under the
null hypothesis of independence $H_0 : X \Indep Y$.
We introduce the corresponding restricted parameter space
$$\CX_\Theta^0 := \{\theta \in \CX_\Theta : \theta_{xy} = \theta_{x+} \theta_{+ y} \ \forall x\in\C{X},y\in\C{Y} \} \subseteq \CX_\Theta.$$
We can also write the null hypothesis as $H_0 : \theta \in \CX_\Theta^0$.
As alternative hypothesis we take that of dependence, i.e., $H_1 : X \nIndep Y$, or equivalently,
$H_1 : \theta \in \CX_\Theta^1$ with $\CX_\Theta^1 := \CX_\Theta \setminus \CX_\Theta^0$.

Suppose now that we have independent and identically distributed data $(X_n,Y_n)_{n=1}^N$ with
$(X_n,Y_n) \given \Theta \sim P(X,Y \given \Theta)$ for all $n=1,\dots,N$, with the ``true'' parameter $\theta$
unknown.
In other words, we assume for the joint distribution on the observed data
$$P((X_n,Y_n)_{n=1}^N \given \Theta) = \bigotimes_{n=1}^N P(X_n,Y_n \given \Theta),$$
where each $P(X_n,Y_n \given \Theta)$ is a copy of the Markov kernel $P(X,Y \given \Theta)$.
We define the \emph{counts} as the number of observations with a given value $(x,y)\in\C{X}\times\C{Y}$:
$$N_{xy} := \sum_{n=1}^N \ind_{(x,y)}(X_n,Y_n).$$
We can represent them in a contingency table:
\begin{center}\begin{tabular}{c|cccc|c}
  & $Y=1$         & $Y=2$ & $\dots$ & $Y=l$ \\
  \hline
  $X=1$ & $N_{11}$      & $N_{12}$      & $\dots$ & $N_{1l}$ & $N_{1+}$\\
  $X=2$ & $N_{21}$      & $N_{22}$      & $\dots$ & $N_{2l}$ & $N_{2+}$\\
  $\vdots$ & $\vdots$     & $\vdots$       & $\ddots$ & $\vdots$  & $\vdots$ \\
  $X=k$ & $N_{k1}$      & $N_{k2}$      & $\dots$ & $N_{kl}$ & $N_{k+}$ \\
  \hline
  & $N_{+ 1}$ & $N_{+ 2}$ & $\dots$ & $N_{+ l}$ & $N_{++} = N$\\
\end{tabular}\end{center}
where we used a similar summation convention for the counts as for the parameters.

The classical frequentist procedure for deciding between two hypotheses $H_0 : \theta \in \CX_\Theta^0$
and $H_1 : \theta \in \CX_\Theta\setminus\CX_\Theta^0$ is as follows.
One comes up with a \emph{test statistic} $T(D)$, which is a function of the data $D \sim P(D \given \Theta)$, 
whose value should help us distinguish between the two hypotheses. 
We will consider here one-sided tests, where a large value of $T(D)$ is in favor of $H_1$ while a small value of $T(D)$ is in favor of $H_0$. 
Then one chooses a particular significance level $\alpha \in (0,1)$. From
the observed data $d$, one then calculates a corresponding $p$-value $p(d)$, which is the probability under 
the null hypothesis that the test statistic has the observed or a more extreme value. For the
one-sided tests we will consider here, the $p$-value can be defined as
$$p(d) := \sup_{\theta\in\CX_\Theta^0} P(T(D) \ge T(d) \given \theta).$$
Then, a decision is taken: if $p(d) \le \alpha$, one considers this as sufficient evidence to reject $H_0$ 
(and accept $H_1$), while if $p(d) > \alpha$, one does not reject $H_0$ as the evidence in the data 
is considered insufficient to do so.
Often, the main desideratum is to control the probability of a Type I error (i.e., the error of incorrectly
rejecting the null hypothesis), which can be achieved by choosing $\alpha$ sufficiently small.
Indeed, from the definition of the $p$-value it follows that:
  \[ \forall \alpha \in (0,1)\, \forall \theta \in \CX_\Theta^0:\;  P( p(D) \le \alpha \given \theta) \le \alpha.\]
For causal discovery, however, we need a more symmetric treatment of the two hypotheses, 
as there we require both the probability of a Type I error and of a Type II error (i.e., the error of incorrectly 
rejecting the alternative hypothesis) to be small. Before we investigate this
tradeoff, let us first propose a concrete test statistic for the case at hand and 
obtain an approximate expression for the corresponding $p$-value.

Here we will work out the details of the \emph{likelihood ratio test}, which for this particular case
is also known as the \emph{$G$ test}. We start by writing down the likelihood:
$$\theta \mapsto P((X_n,Y_n)_{n=1}^N \given \theta) = \prod_{n=1}^N \theta_{X_nY_n} = \prod_{x\in\C{X}} \prod_{y\in\C{Y}} \theta_{xy}^{N_{xy}}$$
where we used the counts as a sufficient statistic of the data. 
This is a multinomial distribution with parameters $(\theta_{xy})_{x\in\C{X},y\in\C{Y}}$ and $N$.
Maximizing the likelihood with respect to the parameters $\theta \in \CX_\Theta$, we obtain the well-known maximum likelihood estimator\footnote{We follow the notational convention in statistics by writing estimators in lowercase (e.g., $\hat \theta$ rather than $\hat \Theta$), even though they are random variables.}
$$\hat \theta_{xy} = \frac{N_{xy}}{N},$$
i.e., the fractions of the different outcomes in the data.
Under the null hypothesis $H_0$, $\theta_{xy} = \theta_{x+}\theta_{+ y}$, and the likelihood factorizes:
\begin{equation*}\begin{split}
  \theta \in \CX_\Theta^0 \implies P((X_n,Y_n)_{n=1}^N \given \theta) 
  &= \prod_{x\in\C{X}} \prod_{y\in\C{Y}} (\theta_{x+} \theta_{+ y})^{N_{xy}} \\
  &= \left(\prod_{x\in\C{X}} \prod_{y\in\C{Y}} \theta_{x+}^{N_{xy}} \right) \left(\prod_{x\in\C{X}} \prod_{y\in\C{Y}} \theta_{+ y}^{N_{xy}}\right) \\
  &= \left(\prod_{x\in\C{X}} \theta_{x+}^{N_{x+}} \right) \left(\prod_{y\in\C{Y}} \theta_{+ y}^{N_{+ y}}\right).
%  &= P_0((X_n)_{n=1}^N \given \theta_{0+},\theta_{1+}) P_0((Y_n)_{n=1}^N \given \theta_{+ 0},\theta_{+ 1})
\end{split}\end{equation*}
This is just the product of two independent multinomial distributions with (variationally independent) parameters
$(\theta_{x+})_{x\in\C{X}}$ and $(\theta_{+ y})_{y\in\C{Y}}$ (and $N$), respectively.
Hence, the restricted maximum likelihood estimator under $H_0$ is 
$$\hat \theta^0_{xy} = \hat \theta^0_{x+} \hat \theta^0_{+ y} = \frac{N_{x+}}{N} \frac{N_{+ y}}{N}.$$
The likelihood ratio is obtained by dividing the likelihood for $\hat \theta$ by the likelihood for $\hat \theta^0$:
\begin{equation}\label{eq:indep_test_XY_likelihood_ratio}\begin{split}
\frac{\sup_{\theta\in\CX_\Theta^0\cup\CX_\Theta^1} P((X_n,Y_n)_{n=1}^N \given \theta)}{\sup_{\theta\in\CX_\Theta^0} P((X_n,Y_n)_{n=1}^N \given \theta)}
  & = \frac{P((X_n,Y_n)_{n=1}^N \given \hat\theta)}{P((X_n,Y_n)_{n=1}^N \given \hat\theta^0)}
    = \prod_{x\in\C{X}} \prod_{y\in\C{Y}} \left(\frac{\hat \theta_{xy}}{\hat \theta^0_{x+} \hat \theta^0_{+ y}}\right)^{N_{xy}} \\
  & = \prod_{x\in\C{X}} \prod_{y\in\C{Y}} \left(\frac{N_{xy} N}{N_{x+} N_{+ y}}\right)^{N_{xy}}.
\end{split}\end{equation}
The likelihood ratio test statistic is defined as 2 times the natural logarithm of this ratio:
\begin{equation}\label{eq:G_test_statistic}
  G_N := 2 \log \frac{\sup_{\theta\in\CX_\Theta^0\cup\CX_\Theta^1} P((X_n,Y_n)_{n=1}^N \given \theta)}{\sup_{\theta\in\CX_\Theta^0} P((X_n,Y_n)_{n=1}^N \given \theta)} = 2 \sum_{x\in\C{X}} \sum_{y\in\C{Y}} N_{xy} \log \frac{N_{xy} N}{N_{x+} N_{+ y}}.
\end{equation}
Since counts can be zero, one should interpret $0 \log \frac{0}{n}$ in this expression as $0$ (for $n \in \mathbb{N}$).

We will now consider the asymptotic behavior of the test statistic under the null hypothesis.
This will yield an approximation for the $p$-value that we can use also for finite samples.
As a simplifying assumption, we will henceforth assume that all probabilities are positive,\footnote{The singularities for vanishing values of $\theta_{xy}$ can be dealt with, but require special attention. For simplicity we study only the regular case here.} i.e.,
\begin{equation}\label{eq:G_test_regularity}
  \theta_{xy} > 0 \qquad \forall x \in \C{X}, y\in \C{Y}.
\end{equation}
\begin{Prp}\label{prp:G_test_asymptotics_H0}
Under $H_0 : X \Indep Y$, and with regularity assumption \eref{eq:G_test_regularity},
$$G_N \leadsto \chi_\nu^2$$
with $\nu = (|\C{X}| - 1)(|\C{Y}| - 1)$ as sample size $N \to \infty$.
  In words, the test statistic $G_N$ converges in distribution\footnote{We say that a sequence of real-valued random variables $X_1,X_2,\dots$ \emph{converges in distribution} to $X_\infty$, and write $X_n \leadsto X_\infty$, if $P(X_n \le x) \to P(X_\infty \le x)$ for all $x \in \RN$ such that $\xi \mapsto P(X_\infty \le \xi)$ is continuous at $x$.} as $N \to \infty$ to a chi-squared distribution with $\nu$ degrees of freedom.\footnote{The chi-square distribution with $\nu$ degrees of freedom is defined as the distribution of a sum of squares of $\nu$ independent standard normal random variables, i.e, of $\sum_{i=1}^\nu Z_i^2$ where $Z_i \sim N(0,1)$ are i.i.d..}
\end{Prp}
\begin{proof}
This is a direct application of Theorem 4.43 in \cite{Bijma++2017}, for which a proof is
provided in Chapter 16 in \cite{Vandervaart1998}.
One has to be careful here to use a different parameterization---in
terms of (variationally) independent parameters---i.e.,
such that the parameter space contains an open part of $\RN^{|\C{X}||\C{Y}|-1}$,
when calculating the score function and the Fisher information matrix
when checking the regularity conditions.
For example, one can choose a pair $(k,l) \in \C{X} \times \C{Y}$ and take
parameters $\theta_{x,y} = \vartheta_{x,y}$ for $x \ne k$ or $y \ne l$, and
$\theta_{k,l} = 1 - \sum_{(x,y) \ne (k,l)} \vartheta_{x,y}$. 
%Then we have an embedding $\RN^{kl-1} \ni \vartheta \mapsto \theta \in \RN^{kl}$.
%Then $\vartheta \in \RN^{kl-1}$.
The dimensionality of $\CX_\Theta$ is $|\C{X}||\C{Y}| - 1$, while
that of $\CX_\Theta^0$ is $(|\C{X}|-1) + (|\C{Y}|-1)$. The degrees
of freedom of the asymptotic chi-square distribution is the difference of the two, i.e.,
$\nu = (|\C{X}||\C{Y}|-1)   -   ((|\C{X}|-1) + (|\C{Y}|-1)) = (|\C{X}| - 1)(|\C{Y}| - 1)$.

An alternative proof will be provided later in a more general setting (see Proposition~\ref{prp:ext_G_test_asymptotics_H0}).
\end{proof}

One therefore obtains an approximate level $\alpha$ test (i.e., a test with Type I error asymptotically upper bounded by $\alpha$) 
by rejecting $H_0$ when $G_N \ge \chi_{\nu,1-\alpha}^2$. Here,
$\chi^2_{\nu,1-\alpha} := F_{\chi^2_\nu}^{-1}(1-\alpha)$ is the upper $\alpha$ quantile of the $\chi^2$-distribution with $\nu$ degrees of freedom, with $F_{\chi^2_\nu}$ the corresponding distribution function (cumulative density function) and $F^{-1}_{\chi^2_\nu}$ its inverse (i.e., the quantile function).
Indeed, if $\theta \in \CX_\Theta^0$, then $P(G_N \ge \chi_{\nu,1-\alpha}^2 \given \theta) \to \alpha$, for any $\alpha \in (0,1)$.
Since 
$$G_N \ge \chi_{\nu,1-\alpha}^2 \iff G_N \ge F_{\chi^2_\nu}^{-1}(1-\alpha) \iff F_{\chi^2_\nu}(G_N) \ge 1-\alpha \iff 1-F_{\chi^2_\nu}(G_N) \le \alpha,$$
the corresponding approximate $p$-value is $1 - F_{\chi^2_\nu}(G_N)$; if this is smaller than or equal to the chosen threshold $\alpha$, 
we reject $H_0$.
This test is called the \emph{$G$-test}.

But what about the Type II error? 
If we let the sample size $N$ grow, we would hope that the probability of a wrong test result becomes arbitrarily small, and
vanishes in the limit $N \to \infty$. 
\begin{Def}\label{def:consistency_ci_test}
A (conditional) independence test is called \emph{consistent} if the probabilities of both Type I and Type II errors converge to 0, no matter what the true parameter value is.
\end{Def}
To obtain consistency, it is not an option to just control Type I error at a fixed level $\alpha$; instead, one has to use a level $\alpha_N$
that depends on the sample size $N$, and converges to 0 (implying that Type I error converges to 0). However, because of the
tradeoff between Type I and Type II errors, the rate at which $\alpha_N$ converges to $0$ has to be chosen carefully in order to
be able to guarantee that also Type II error vanishes asymptotically. As we shall see, the convergence rate of $\alpha_N$ should be chosen
sufficiently slow.

While it is often easier to calculate the Type I error than the Type II error of
a test, in this case we can actually analyze the asymptotic behavior of the test statistic under the alternative hypothesis $H_1$.
Define
$$\hat I_N := \sum_{x\in\C{X}} \sum_{y\in\C{Y}} \hat \theta_{xy} \log \frac{\hat \theta_{xy}}{\hat \theta_{x+} \hat \theta_{+ y}}
%    = \sum_{x\in\C{X}} \sum_{y\in\C{Y}} \frac{N_{xy}}{N} \log \frac{N_{xy} N}{N_{x +} N_{+ y}}
    = \frac{G_N}{2N}$$
where we used that $\hat\theta^0_{x +} = \hat\theta_{x +}$ and $\hat\theta^0_{+ y} = \hat\theta_{+ y}$.
This is an estimator (the so-called ``plug-in estimator'' $I(\hat\theta)$) of the mutual information $I(X;Y)$:
\begin{equation*}\begin{split}
  I(\theta) & := \sum_{x\in\C{X}}\sum_{y\in\C{Y}} \theta_{xy} \log \frac{\theta_{xy}}{\theta_{x +} \theta_{+ y}} \\
  & = \sum_{x\in\C{X}} \sum_{y\in\C{Y}} P(X=x,Y=y \given \theta) \log \frac{P(X=x,Y=y \given \theta)}{P(X=x \given \theta) P(Y=y \given \theta)} =: I(X;Y \kgiven \theta).
\end{split}\end{equation*}
With Jensen's inequality, one can show that $I(X;Y \kgiven \theta) \ge 0$, and that $I(X;Y \kgiven \theta) = 0 \iff X \Indep Y \given \Theta=\theta$.
Note further that the function $\CX_\Theta \to [0,\infty) : \theta \mapsto I(\theta)$ is continuous.

With this observation, we can prove the asymptotic consistency of the $G$-test under assumptions on the critical values used
for deciding between $H_0$ and $H_1$. 
\begin{Cor}\label{cor:G_test_consistency}
Consider an infinite sequence of $G$ tests performed on the first $N$ samples of an infinitely large
data set $(X_n,Y_n)_{n=1}^\infty$, where one 
accepts $H_1 : X \nIndep Y$ if $G_N \ge \tau_N$, and otherwise accepts $H_0 : X \Indep Y$, 
for some given sequence of thresholds $\tau_N$.
Under the regularity assumption \eref{eq:G_test_regularity}, this sequence of tests is
asymptotically consistent if $\tau_N \to \infty$ but $\tau_N / N \to 0$.
\end{Cor}
\begin{proof}
We start by a simple application of the strong law of large numbers.
Let $\theta \in \CX_\Theta$. Since the $(X_n,Y_n)$ are assumed to be i.i.d., and 
  $$\Exp \big(\I_{(x,y)}(X_n,Y_n)\big) = \theta_{xy}$$
for all $x \in \C{X}, y \in \C{Y}$,
we conclude that $N_{xy} / N \stackrel{a.s.}{\to} \theta_{xy}$ for all $x \in \C{X}, y \in \C{Y}$
by the strong law of large numbers.\footnote{The convergence is ``almost surely'', 
i.e., $N_{xy} / N \stackrel{a.s.}{\to} \theta_{xy}$ means that
$P(N_{xy} / N \to \theta_{xy} \given \theta) = 1$.}
Hence $\hat\theta_{xy} \stackrel{a.s.}{\to} \theta_{xy}$. Hence, also
$\hat\theta_{+ y} \stackrel{a.s.}{\to} \theta_{+ y}$ and
$\hat\theta_{x +} \stackrel{a.s.}{\to} \theta_{x +}$.
Furthermore, by continuity, $I(\hat\theta) \stackrel{a.s.}{\to} I(\theta)$.
Hence, $G_N / N \stackrel{a.s.}{\to} 2I(\theta)$. 
  
Under $H_1$, we have $I(\theta) > 0$,
and since by assumption $\tau_N / N \to 0$, 
$\I_{G_N < \tau_N} \stackrel{a.s.}{\to} 0$. 
Since a.s.\ convergence implies convergence in probability,
  $$\theta \in \CX_\Theta^1 \implies P(G_N < \tau_N \given \theta) \to 0.$$
Thus, the probability of a Type II error vanishes asymptotically.

This same approach doesn't work for the Type I error. The reason is that even though
we assume $\tau_N \to \infty$, and we know that $G_N / N \stackrel{a.s.}{\to} 0$ 
under $H_0$,
this does not suffice to conclude anything about the probability of the event 
$G_N \ge \tau_N$. But we can make use of Proposition~\ref{prp:G_test_asymptotics_H0},
which states that $G_N \leadsto \chi_{\nu}^2$ under $H_0$. 
Since the distribution function of $\chi_{\nu}^2$ is continuous,
this implies uniform convergence of the distribution functions:
  $$\sup_{x \in \RN} |F_{G_N}(x) - F_{\chi_{\nu}^2}(x)| \to 0.$$
Hence
  $$|F_{G_N}(\tau_N) - F_{\chi_{\nu}^2}(\tau_N)| \le \sup_{x\in\RN} |F_{G_N}(x) - F_{\chi_{\nu}^2}(x)| \to 0.$$
Since $\tau_N \to \infty$, $F_{\chi_{\nu}^2}(\tau_N) \to 1$.
Hence, also $F_{G_N}(\tau_N) \to 1$.
%Since $\tau_N \to F^{-1}_{\chi_{\nu}^2}(1-\alpha)$ and $F_{\chi_{\nu}^2}$ is continuous,
%  $$F_{\chi_\nu^2}(\tau_N) \to F_{\chi_\nu^2}(F^{-1}_{\chi_\nu^2}(1-\alpha)) = 1-\alpha.$$
%Combining these two, we conclude
%  $$F_{G_N}(\tau_N) \to 1-\alpha$$
%and hence the probability of a Type I error converges to $\alpha$,
%  $$\Prb(G_N > \tau_N) \to \alpha.$$
We conclude that 
  $$\theta \in \CX_\Theta^0 \implies P(G_N \ge \tau_N \given \theta) \to 0,$$
i.e., the probability of a Type I error converges to 0.
%Under $H_1$, Proposition~\ref{prp:G_test_asymptotics_H1} states that
%  $$\sqrt{N}(\hat I_N - I) \leadsto N(0,\sigma^2).$$
%Again, since the normal distribution is continuous, this implies uniform convergence
%  $$\sup_{x\in\RN} |F_{\sqrt{N}(\hat I_N - I)} - F_{N(0,\sigma^2)}| \to 0.$$
%Since $\hat I_N = \frac{G_N}{2N}$, this implies
%  $$\sup_{x\in\RN} |F_{\hat G_N} - F_{N(2NI,4N\sigma^2)}| \to 0.$$
%Hence,
%  $$|F_{\hat G_N}(\tau_N) - F_{N(2NI,4N\sigma^2)}(\tau_N)| \to 0.$$
%As we assumed that $\tau_N / N \to 0$, $F_{N(2NI,4N\sigma^2)}(\tau_N) \to 0$
%and hence the asymptotic probability of a Type II error (incorrectly rejecting $H_1$) converges to 0,
%  $$\Prb(\hat G_N \le \tau_N) \to 0.$$
\end{proof}

While one traditionally focuses mostly on Type I error control, in causal discovery we
are more interested in having both small Type I and Type II error. In order to achieve this
(at least asymptotically, i.e., for sufficiently large sample sizes), we can thus make use of a sequence of
thresholds that satisfies the assumptions in the corollary. In terms of $p$-values, this means
that to bound the Type I error, a fixed critical value $\alpha$ suffices, but for consistency we
let $\alpha_N \to 0$ with a rate such that $\chi^2_{\nu,1-\alpha_N} / N \to 0$.

While for a finite sample, we can give guarantees (at least approximately) on the Type I error, it will often be
impossible to provide guarantees on the Type II error without making strong
assumptions on the parameters. Indeed, since the mutual information $I(X;Y)$ (a measure of
the dependence of $X$ and $Y$) can be arbitrarily close to zero for weakly dependent $X$ 
and $Y$, one cannot know in advance how many samples will be needed to be able to distinguish
it from an independence.\footnote{This is referred to as the lack of ``uniformly consistent''
(conditional) independence tests.}

\subsection{Conditional Independence for Categorical Random Variables}

We now extend the $G$ test to a conditional independence test that we will refer to as the conditional $G$ test.

Consider three categorical random variables $X,Y,Z$ taking values in spaces $\C{X}$, $\C{Y}$
and $\C{Z}$, respectively (with $2 \le |\C{X}| < \infty$, $2 \le |\C{Y}| < \infty$ and $1 \le |\C{Z}| < \infty$) and joint distribution $P(X,Y,Z)$.
With Remark~\ref{ci-random-variables}, we get (because finite spaces are standard):
\begin{equation*}\begin{split}
  X \Indep_{P(X,Y,Z)} Y \given Z 
  & \iff P(X,Y,Z) = P(X | Z) \otimes P(Y,Z) \\
  & \iff P(X,Y | Z) = P(X | Z) \otimes P(Y | Z) \quad \text{$P(Z)$-a.s.} \\
  & \begin{split}
    {} \iff \forall z \in \C{Z}: \big[ & P(Z=z) > 0 \implies \\
    & P(X,Y \given Z=z) = P(X \given Z=z) P(Y \given Z=z) \big]
  \end{split} \\
  & \iff \forall z \in \C{Z}: \big[P(Z=z) > 0 \implies X \Indep_{P(X,Y \given Z=z)} Y \big].
\end{split}\end{equation*}
This suggests that we can make use of an independence test for two categorical variables
on each  ``stratum'' corresponding to conditioning on a specific value $Z=z$ that has
positive probability to occur.

We parameterize the conditional kernel $P(X,Y | Z)$ in terms of parameters
$(\theta_{xy|z})_{x\in\C{X},y\in\C{Y},z\in\C{Z}}$ which live in space
$$\CX_{\Theta_{XY|Z}} := %\{ \theta \in \prod_{z\in\C{Z}} [0,1] : \sum_{z\in\C{Z}} \theta_z = 1\} \times
\{ \theta \in \prod_{x\in\C{X},y\in\C{Y},z\in\C{Z}} [0,1] : \forall_{z\in\C{Z}} \sum_{x\in\C{X},y\in\C{Y}} \theta_{xy|z} = 1\}.$$
With the summation convention, we can write the normalization condition as $\theta_{++|z} = 1$ for all $z \in \C{Z}$.
For those $z\in\C{Z}$ with $P(Z = z) > 0$, we have 
$$\frac{P(X = x, Y = y, Z = z \given \Theta = \theta)}{P(Z=z \given \Theta=\theta)} = P(X = x, Y = y \given Z=z, \Theta=\theta) = \theta_{xy|z}.$$
We also parameterize the marginal distribution $P(Z)$ in terms of parameters
$(\theta_z)_{z\in\C{Z}}$ which live in space
$$\CX_{\Theta_Z} := \{ \theta \in \prod_{z\in\C{Z}} [0,1] : \sum_{z\in\C{Z}} \theta_z = 1\}.$$
Any joint distribution of $X,Y$ and $Z$ can then be parameterized as
$$P(X=x,Y=y,Z=z \given \Theta=\theta) = \theta_z \theta_{xy|z},$$
with parameter space
$$\CX_\Theta := \CX_{\Theta_Z}\times \CX_{\Theta_{XY|Z}}.$$

We formulate the null hypothesis $H_0 : X \Indep Y \given Z$ of independence in terms of the parameters as
$$\forall x \in \C{X}, y \in \C{Y}, \forall z \in \C{Z}: \quad \theta_{xy|z} = \theta_{x+|z} \theta_{+ y|z}$$
(for convenience, we have strengthened it a bit; strictly speaking, we only need this relation to hold
for all $z \in \C{Z}$ with $\theta_z > 0$; however, since the data will not convey any information on $\theta_{xy|z}$ for such $z$, this does not matter).
The corresponding restricted parameter space is
$$\CX_{\Theta_{XY|Z}}^0 := \{\theta \in \CX_\Theta : \theta_{xy|z} = \theta_{x+|z} \theta_{+ y|z}\ \forall x \in \C{X}, y \in \C{Y}, z \in \C{Z}\}.$$
We can then also write the null hypothesis as $H_0 : \theta \in \CX_\Theta^0$ with
$\CX_\Theta^0 := \CX_{\Theta_Z} \times \CX_{\Theta_{XY\given Z}}^0$.
As alternative hypothesis we take that of dependence, i.e., $H_1 : X \nIndep Y \given Z$, or equivalently,
$H_1 : \theta \in \CX_\Theta^1$, where
$\CX_\Theta^1 := \CX_{\Theta_Z} \times \CX_{\Theta_{XY\given Z}}^1$ 
with $\CX_{\Theta_{XY\given Z}}^1 := \CX_{\Theta_{XY\given Z}} \setminus \CX_{\Theta_{XY\given Z}}^0$.

Suppose now that we have independent and identically distributed data $(X_n,Y_n,Z_n)_{n=1}^N$ with
$(X_n,Y_n,Z_n) \sim P(X,Y,Z \given \Theta)$ for all $n=1,\dots,N$, with the ``true'' parameter $\theta \in \CX_\Theta$ unknown.
In other words, we assume for the joint distribution on the observed data
$$P((X_n,Y_n,Z_n)_{n=1}^N \given \Theta) = \bigotimes_{n=1}^N P(X_n,Y_n,Z_n \given \Theta),$$
where each $P(X_n,Y_n,Z_n \given \Theta)$ is a copy of the Markov kernel $P(X,Y,Z \given \Theta)$.
We define the \emph{counts} as the number of observations with a given value $(x,y,z)\in\C{X}\times\C{Y}\times\C{Z}$:
$$N_{xyz} := \sum_{n=1}^N \ind_{(x,y,z)}(X_n,Y_n,Z_n).$$

We again work out the details of the likelihood ratio test, and start by writing down the likelihood:
\begin{equation*}\begin{split}
  \theta \mapsto P((X_n,Y_n,Z_n)_{n=1}^N \given \theta) 
  &= \prod_{n=1}^N (\theta_{X_n,Y_n|Z_n} \theta_{Z_n}) 
   = \prod_{x\in\C{X}} \prod_{y\in\C{Y}} \prod_{z\in\C{Z}} (\theta_{xy|z}^{N_{xyz}} \theta_z^{N_{xyz}}) \\
  &= \left( \prod_{z \in \C{Z}} \theta_z^{N_{+ + z}} \right) \left(\prod_{z \in \C{Z}}\prod_{x\in\C{X}} \prod_{y\in\C{Y}} \theta_{xy|z}^{N_{xyz}}\right)
\end{split}\end{equation*}
where we used the counts as a sufficient statistic of the data. 
We recognize the first factor as the likelihood of a multinomial
distribution with parameters $(\theta_z)_{z\in\C{Z}}$ and $N$.
The second factor is a product of the likelihoods
of multinomial distributions with parameters $(\theta_{xy|z})_{x\in\C{X},y\in\C{Y}}$ and $N_{++ z}$,
for each $z \in \C{Z}$.
Maximizing the likelihood with respect to the parameters $\theta \in \CX_\Theta$, we obtain the maximum likelihood estimator
$$\left(\hat \theta_{xy|z},\hat \theta_z\right) = \left(\frac{N_{xyz}}{N_{++ z}},\frac{N_{++ z}}{N}\right).$$
Under the null hypothesis $H_0$, $\theta_{xy|z} = \theta_{x+|z}\theta_{+ y|z}$, and the likelihood factorizes over $X$
and $Y$:
\begin{equation*}\begin{split}
  \theta \in \CX_\Theta^0 \implies P((X_n,Y_n,Z_n)_{n=1}^N \given \theta) 
  &= \left( \prod_{z \in \C{Z}} \theta_z^{N_{+ + z}} \right) \prod_{z \in \C{Z}}\prod_{x\in\C{X}} \prod_{y\in\C{Y}} (\theta_{x+|z} \theta_{+ y|z})^{N_{xyz}} \\
  &= \left( \prod_{z \in \C{Z}} \theta_z^{N_{+ + z}} \right) \prod_{z \in \C{Z}} \left(\prod_{x\in\C{X}} \theta_{x+|z}^{N_{x+ z}} \right) \left(\prod_{y\in\C{Y}} \theta_{+ y|z}^{N_{+ yz}} \right).
\end{split}\end{equation*}
The restricted maximum likelihood estimator under $H_0$ is 
$$\left(\hat \theta^0_{xy|z},\hat \theta^0_z\right) = \left(\hat \theta^0_{x+|z} \hat \theta^0_{+ y|z},\hat \theta^0_z\right) = \left(\frac{N_{x+ z} N_{+ y z}}{N_{++ z}^2},\frac{N_{++ z}}{N}\right).$$
The likelihood ratio is obtained by dividing the likelihood for $\hat \theta$ by the likelihood for $\hat \theta^0$:
\begin{equation*}\begin{split}
\frac{\sup_{\theta\in\CX_\Theta^0\cup\CX_\Theta^1} P((X_n,Y_n,Z_n)_{n=1}^N \given \theta)}{\sup_{\theta\in\CX_\Theta^0} P((X_n,Y_n,Z_n)_{n=1}^N \given \theta)}
  & = \frac{P((X_n,Y_n,Z_n)_{n=1}^N \given \hat\theta)}{P((X_n,Y_n,Z_n)_{n=1}^N \given \hat\theta^0)}
    = \prod_{z\in\C{Z}} \prod_{x\in\C{X}} \prod_{y\in\C{Y}} \left(\frac{\hat \theta_{xy|z}}{\hat \theta^0_{x+|z} \hat \theta^0_{+ y|z}}\right)^{N_{xyz}} \\
  & = \prod_{z\in\C{Z}} \prod_{x\in\C{X}} \prod_{y\in\C{Y}} \left(\frac{N_{xyz} N_{++ z}}{N_{x+ z} N_{+ y z}}\right)^{N_{xyz}},
\end{split}\end{equation*}
where the factors involving the marginal $P(Z)$ cancel out.
The likelihood ratio test statistic is defined as 2 times the natural logarithm of this ratio:
\begin{equation}\label{eq:cond_G_test_statistic}
G_N := 2 \log \frac{\sup_{\theta\in\CX\Theta^0\cup\CX_\Theta^1} P((X_n,Y_n,Z_n)_{n=1}^N \given \theta)}{\sup_{\theta\in\CX_\Theta^0} P((X_n,Y_n,Z_n)_{n=1}^N \given \theta)} = 2 \sum_{z\in\C{Z}} \sum_{x\in\C{X}} \sum_{y\in\C{Y}} N_{xyz} \log \frac{N_{xyz} N_{+ + z}}{N_{x+ z} N_{+ y z}}.
\end{equation}

We will now consider the asymptotic behavior of the test statistic under both hypotheses.
As a simplifying assumption, we will henceforth assume that all probabilities are positive, i.e.,
\begin{equation}\label{eq:cond_G_test_regularity}
  \begin{cases}
    \theta_z > 0 & \forall z \in \C{Z}, \text{ and}\\
    \theta_{xy|z} > 0 & \forall x \in \C{X}, y\in \C{Y}, z\in \C{Z}.
  \end{cases}
\end{equation}
\begin{Prp}\label{prp:cond_G_test_asymptotics_H0}
Under $H_0 : X \Indep Y \given Z$, and with regularity assumption \eref{eq:cond_G_test_regularity}
$$G_N \leadsto \chi_\nu^2$$
with $\nu = |\C{Z}|(|\C{X}| - 1)(|\C{Y}| - 1)$ as sample size $N \to \infty$.
In words, the test statistic $G_N$ converges in distribution as $N \to \infty$ to a chi-squared distribution with $\nu$ degrees of freedom.
\end{Prp}
\begin{proof}
This is analogous to the proof of Proposition~\ref{prp:G_test_asymptotics_H0}.
The dimensionality of $\CX_\Theta$ is $|\C{Z}|(|\C{X}||\C{Y}| - 1) + (|\C{Z}|-1)$, while
that of $\CX_\Theta^0$ is $|\C{Z}|\big((|\C{X}|-1) + (|\C{Y}|-1)\big) + (|\C{Z}|-1)$. The degrees
of freedom of the asymptotic chi-square distribution is the difference of the two, i.e.,
$\nu = |\C{Z}|(|\C{X}| - 1)(|\C{Y}| - 1)$.
\end{proof}

One therefore obtains an approximate level $\alpha$ test (i.e., a test with Type I error asymptotically upper bounded by $\alpha$) 
by rejecting $H_0$ when $G_N \ge \chi_{\nu,1-\alpha}^2$.

Define
$$\hat I_N := \sum_{z\in\C{Z}} \hat\theta_{z} \sum_{x\in\C{X}} \sum_{y\in\C{Y}} \hat \theta_{xy|z} \log \frac{\hat \theta_{xy|z}}{\hat \theta_{x+|z} \hat \theta_{+ y|z}}
%    = \sum_{x\in\C{X}} \sum_{y\in\C{Y}} \frac{N_{xy}}{N} \log \frac{N_{xy} N}{N_{x +} N_{+ y}}
    = \frac{G_N}{2N}$$
where we used that $\hat\theta^0_{x +|z} = \hat\theta_{x +|z}$ and $\hat\theta^0_{+ y|z} = \hat\theta_{+ y|z}$.
This is a plug-in estimator of the conditional mutual information $I(X;Y|Z)$:
\begin{equation*}\begin{split}
  I(\theta) 
  & := \sum_{z\in\C{Z}} \theta_z \sum_{x\in\C{X}}\sum_{y\in\C{Y}} \theta_{xy|z} \log \frac{\theta_{xy|z}}{\theta_{x +|z} \theta_{+ y|z}} \\
  & = \sum_{z\in\C{Z}} \sum_{x\in\C{X}} \sum_{y\in\C{Y}} P(X=x,Y=y,Z=z \given \Theta=\theta) \log \frac{P(X=x,Y=y|Z=z,\Theta=\theta)}{P(X=x|Z=z,\Theta=\theta) P(Y=y|Z=z,\Theta=\theta)} \\
  & =: I(X;Y|Z\kgiven \theta).
\end{split}\end{equation*}
With Jensen's inequality, one can show that $I(X;Y|Z) \ge 0$, and that $I(X;Y|Z) = 0 \iff X \Indep Y \given Z$.
Note further that the function $\CX_\Theta \to [0,\infty) : \theta \mapsto I(\theta)$ is continuous.

With this observation, we can prove the asymptotic consistency of the conditional $G$-test under assumptions on the critical values used
for deciding between $H_0$ and $H_1$.
\begin{Cor}\label{cor:cond_G_test_consistency}
Consider an infinite sequence of conditional $G$ tests performed on the first $N$ samples of an infinitely large
data set $(X_n,Y_n,Z_n)_{n=1}^\infty$,
where one accepts $H_1 : X \nIndep Y \given Z$ if $G_N \ge \tau_N$, and otherwise accepts $H_0 : X \Indep Y \given Z$, 
for some given sequence of thresholds $\tau_N$.
Under the regularity assumption \eref{eq:cond_G_test_regularity}, this sequence is
asymptotically consistent if $\tau_N \to \infty$ but $\tau_N / N \to 0$.
\end{Cor}
\begin{proof}
This is very similar to the proof of Corollary~\ref{cor:G_test_consistency}.
  
We again apply the strong law of large numbers.
Let $\theta \in \CX_\Theta$. Since $(X_n,Y_n,Z_n)$ are assumed to be i.i.d., and 
  $$\Exp \big(\I_{(x,y,z)}(X_n,Y_n,Z_n)\big) = \theta_{xy|z}\theta_z$$
for all $x \in \C{X}, y \in \C{Y}, z \in \C{Z}$,
we conclude that $N_{xyz} / N \stackrel{a.s.}{\to} \theta_{xy|z} \theta_z$ for all $x \in \C{X}, y \in \C{Y}, z\in \C{Z}$
by the strong law of large numbers. Hence, also $N_{+ + z} / N \stackrel{a.s.}{\to} \theta_z$ for all $z \in \C{Z}$.
Hence, using \eref{eq:cond_G_test_regularity},
$\hat\theta_z \stackrel{a.s.}{\to} \theta_z$,
$\hat\theta_{xy|z} \stackrel{a.s.}{\to} \theta_{xy|z}$,
$\hat\theta_{+ y|z} \stackrel{a.s.}{\to} \theta_{+ y|z}$, and
$\hat\theta_{x +|z} \stackrel{a.s.}{\to} \theta_{x +|z}$.
By continuity, $I(\hat\theta) \stackrel{a.s.}{\to} I(\theta)$.
Hence, $G_N / N \stackrel{a.s.}{\to} 2I(\theta)$. 
  
We can now reason analogously as in the proof of 
Corollary~\ref{cor:G_test_consistency} to conclude that
%Under $H_1$, we have $I(\theta) > 0$,
%and since by assumption $\tau_N / N \to 0$, 
%$\I_{G_N \le \tau_N} \stackrel{a.s.}{\to} 0$. 
%Since a.s.\ convergence implies convergence in probability,
%  $$\Prb_{H_1}(G_N \le \tau_N) \to 0.$$
%Thus, 
the probability of a Type II error vanishes asymptotically.

For an asymptotic estimate of the probability of a Type I error,
we can make use of Proposition~\ref{prp:cond_G_test_asymptotics_H0},
which states that $G_N \leadsto \chi_{\nu}^2$. 
This part of the proof is identical to the corresponding part of the proof of Corollary~\ref{cor:G_test_consistency}.
\end{proof}
